\section{Preliminaries}

Most instances of the problems considered in this work consist of a graph $G=\br{V,E,D,c,w}$, where $V$ is the vertex set, $E$ is the set of supply edges, $D\subseteq\binom{V}{2}$ is an undirected loopless demand graph, $c\colon E \to \mathbb{N}$ is the \emph{capacity} function which encodes the cost of cutting a supply edge, and $w\colon D\to \mathbb{N}$ assigns a nonnegative weight to each demand edge. Every demand edge is counted once.
For adjacent vertices $u,v\in V$, we denote by $uv$ the supply edge between $u$ and $v$.
For any $C\subseteq E$, $G\setminus C$ is the graph obtained by removing the edges in $C$ from $G$.
We write $\br{V_H,E_H,D_H,c_H,w_H}$ to point out that given parameters represent a particular graph or subgraph $H$, different than $G$.
For $U\subseteq V$, the induced subgraph $G[U]$ has supply-edge set $E[U]=\brc{uv\in E:u,v\in U}$, demand-edge set $D[U]=D\cap\binom{U}{2}$, and the restrictions of $c$ and $w$ to these sets.

For an edge set $E'\subseteq\binom{V}{2}$ and a vertex set $S\subseteq V$, let
\[
	\delta_{E'}(S)
	=\brc{uv\in E':\spr{\brc{u,v}\cap S}=1}
\]
be the edges of $E'$ with exactly one endpoint in $S$. For a partition $\cP$ of $V$, let
\[
	\delta_{E'}(\cP)
	=\brc{uv\in E':u\text{ and }v\text{ belong to distinct parts of }\cP}.
\]
For $A\subseteq V$, let
\[
	D\br{A}
	=\brc{uv\in D:\brc{u,v}\cap A\neq\emptyset}
	=D[A]\cup\delta_D(A)
\]
denote the demand edges incident to $A$. Thus $D[A]$ contains only demand edges internal to $A$, whereas $D\br{A}$ also contains the demand edges crossing the cut $\br{A,\bar A}$.

For disjoint subsets $A, B\subseteq V$, we define
\[
w\br{A}=\sum_{\substack{uv\in D\\u,v\in A}}w\br{u,v},
\qquad
w\br{A,B}=\sum_{\substack{uv\in D\\u\in A,\ v\in B}}w\br{u,v}.
\]
Whenever $A, B\subseteq V$ are not disjoint, we define $w\br{A,B}=w\br{A}+w\br{A,B\setminus A}$ (see Figure~\ref{fig:demand-decomposition} for an illustration).
For a vertex $v\in V$, define its weighted demand degree as
\[
d_w(v)=w\br{\delta_D\br{\brc{v}}},
\]
and, for $S\subseteq V$, put $d_w(S)=\sum_{v\in S}d_w(v)$. Since every internal demand edge has both endpoints in $S$ and every crossing demand edge has one endpoint in $S$, we have
\begin{equation}\label{eq:demand-degree-identity}
	d_w(S)
	=2w\br{S}+w\br{S,V\setminus S}
	=w\br{V}+w\br{S}-w\br{V\setminus S}.
\end{equation}
For a subset $C\subseteq E$ and a capacity function $c$ on the supply edges, $c\br{C}=\sum_{e\in C}c\br{e}$.
For a subset $S\subseteq V$, we denote by $\bar{S}$ the complement of $S$ in $V$, i.~e. $\bar{S}=V\setminus S$.
A \emph{cut} is any subset $S\subseteq V$, such that $S\neq \emptyset$, $S\neq V$ and we denote it either by just $S$ or we write $\br{S, \bar{S}}$ to explicitly indicate the two sides of the cut.
It will be sometimes more convenient to call the edge set $\delta_E(S)$ a cut.
Which definition is used will be clear from the context.
For any cut-based objective function $F$ used in this paper, $F\br{S}$ or $F\br{S, \bar{S}}$ denote the value of $F$ on a specific cut $S$, whereas $F\br{G}$ denotes the optimum value of $F$ over all cuts of the graph $G$.
Moreover, whenever an objective function is given as a ratio and its denominator equals $0$, we set the value of this objective to $\infty$.

\begin{figure}[H]
  \centering
  \begin{tikzpicture}[>=stealth,line width=0.9pt,
      vnA/.style={circle,draw=cutA,fill=cutA!12,minimum size=4.5mm,inner sep=0pt,line width=0.8pt,
          preaction={fill=black,opacity=0.32,transform canvas={xshift=1.3pt,yshift=-1.3pt}}},
      vnB/.style={circle,draw=cutB,fill=cutB!12,minimum size=4.5mm,inner sep=0pt,line width=0.8pt,
          preaction={fill=black,opacity=0.32,transform canvas={xshift=1.3pt,yshift=-1.3pt}}}]
        \draw[draw=cutA,fill=cutA!8,line width=1.2pt,
          preaction={fill=black,opacity=0.22,transform canvas={xshift=1.5pt,yshift=-1.5pt}}]
      (-3.1,0) ellipse (2.1 and 1.4);
        \draw[draw=cutB,fill=cutB!8,line width=1.2pt,
          preaction={fill=black,opacity=0.22,transform canvas={xshift=1.5pt,yshift=-1.5pt}}]
      (3.3,0) ellipse (2.1 and 1.4);
    \node[font=\bfseries,text=cutA] at (-3.1,1.7) {$S$};
    \node[font=\bfseries,text=cutB] at (3.3,1.7) {$\bar S$};

    \node[vnA] (s1) at (-3.7, 0.7) {};
    \node[vnA] (s2) at (-2.9, 0.8) {};
    \node[vnA] (s3) at (-2.3, 0.0) {};
    \node[vnA] (s4) at (-3.0,-0.7) {};
    \draw[cutA!70] (s1)--(s2)--(s3)--(s4)--(s1) (s1)--(s3) (s2)--(s4);

    \node[vnB] (t1) at (2.3, 0.7) {};
    \node[vnB] (t2) at (3.3, 0.8) {};
    \node[vnB] (t3) at (3.9, 0.1) {};
    \node[vnB] (t4) at (3.3,-0.7) {};
    \node[vnB] (t5) at (2.4,-0.7) {};
    \draw[cutB!70] (t1)--(t2)--(t3)--(t4)--(t5)--(t1) (t1)--(t3) (t2)--(t4) (t3)--(t5);

    \draw[cutC,dashed,line width=1.1pt] (s2) -- (t1);
    \draw[cutC,dashed,line width=1.1pt] (s2) to[out=-12,in=160] (t4);
    \draw[cutC,dashed,line width=1.1pt] (s3) to[out=-42,in=222] (t2);
    \draw[cutC,dashed,line width=1.1pt] (s4) to[out=-18,in=200] (t5);
  \end{tikzpicture}
  \caption{$f(S,V)=\textcolor{cutA}{f(S)}+\textcolor{cutC}{f(S,\bar S)}$: the total `mass' incident to $S$ decomposes into the mass inside $S$ and the mass crossing the cut.}
  \label{fig:demand-decomposition}
\end{figure}

